%!TEX program = lualatex
\documentclass[ngerman,headheight=70pt]{scrartcl}
\usepackage[ngerman]{babel}
\usepackage{blindtext}
\usepackage[utf8]{luainputenc}
\usepackage{scrlayer-scrpage}
\usepackage[T1]{fontenc}
\usepackage{fontspec}
\usepackage{multicol}
\usepackage{enumitem}
\usepackage[german=quotes]{csquotes}
\usepackage[bottom=2cm,footskip=8mm]{geometry}
\setmainfont{Source Sans Pro}
\MakeOuterQuote{"}

\setlength{\parindent}{0pt}
\setlength{\parskip}{6pt}

\pagestyle{plain.scrheadings}
\thispagestyle{scrheadings}
\lohead{\vspace*{-1.5cm}\includegraphics[scale=0.4]{../images/stupa_logo_rot.pdf}}
\rohead{\\Präsidium des\\Studierendenparlaments}

\lofoot[Universität Hamburg · Präsidium des Studierendenparlaments\\\hrule Von-Melle-Park 5 · D-20146 Hamburg]{Universität Hamburg · Präsidium des Studierendenparlaments\\\hrule Von-Melle-Park 5 · D-20146 Hamburg}
\cofoot[]{}
\cohead[\thepage]{}

\renewcommand{\thesubsection}{TOP \arabic{subsection}}
\renewcommand{\thesubsubsection}{\arabic{subsubsection}.}
\setcounter{subsection}{-1}

\begin{document}
    UHH · StuPa-Präsidium · Von-Melle-Park 5 · D-20146 Hamburg

    %%% PAGE ONE (EDIT LATER) %%%%
    % Hier den Titel des Protokolls eintragen
    \section*{Protokoll der Konstituierenden Sitzung des Studierendenparlaments vom 06. April 2017}

    % Protokollanten eintragen
    \textbf{Protokoll: Gunhild Berdal, Ramon Weilinger}\\
    \textbf{Ort: VMP9 Hörsaal}\\
    \textbf{Beginn: 18. Uhr}\\
    \textbf{Ende: 23. Uhr}

    % Wohl die komplizierteste Stelle im Protokoll
    % Als Faustregel sollten nicht mehr als 3 Namen auf einer Zeile stehen
    % Einfach an dieser Vorlage orientieren
    % Im Idealfall ist hier nicht viel zu ändern, da kaum jemand fehlt.
    % Rücktritte müssen natürlich pro Protokoll angepasst werden
    \vspace{0.5cm}
    \begin{tabular}{ll}
        Anwesend: & \\
            RCDS (5 Sitze): & Ramon Weilinger \\
             CampusGrün (14 Sitze): & Laura Franzen, Geoffrey Youett, Elena Rysikova, \\
                                   & Philipp Droll, Yasemin Günther, Melf Johannsen, \\
                                   & Sonja Lattwesen, Mario Moldenhauer, Christian Lagod, \\
                                   & Jim Martens, Svenja Horn, Damian Hofmann, \\
                                   & Martin Sievert  \\
             BGZM-Liste (2 Sitze) : & \\
             WiWi (2 Sitze): & \\
             Unicorns (5 Sitze): &  \\
             Liste LINKS (3 Sitze): & Gunhild Berdal, Till Petersen, Sinah Mielich \\
             HWP (2 Sitze): &  \\
             MIN (4 Sitze): &  \\
             SDS* (3 Sitze): & Mena Winkler, Jacob Petersein, Artur Brückmann \\
             Bart-LISTE (2 Sitze): & Timo Zeimet, Malte Peters \\
             LHG (1 Sitz): & Tobias Heisig \\
             harte zeiten (1 Sitz): & Tobias Berking \\
             Jura (1 Sitz): & \\
             AL (2 Sitze): & Henri Weber \\
            & \\
        Entschuldigt: & Jennifer Maack (RCDS), Armin Günther (CG), Jakob Pape, \\
                      & Karen Martirosian (BGZM), Johannes Peplow, Andreas Hartkamp (Unicorns), \\
                      & Jan Detampel (MIN), Karima Schulze (AL)\\
                      &\\
        Unentschuldigt abwesend: & Kolja Kolb, Ramin Shakiba, Benjamin Welling (RCDS), \\
                                & Claas-Friso Hente, Charleen Lorenz (WiWi), Katharina Kucza, \\
                                & Annkathrin Löffler, Marielle Hermstrüwer (Unicorns), Hakan Tasdemir, \\
                                & Gökce Celik (HWP), Ailina Salten, Lotte Rullkötter, \\
                                & Nabil Jabareen (MIN), Maurice Martens (Jura)\\
        Rücktritte: & \\
        Mandatsverluste: & \\
    \end{tabular}
    \newpage
    %%% Vorgeschlagene Tagesordnung %%%
    \underline{Vorgeschlagene Tagesordnung}
    \begin{enumerate}[label={\textbf{Top \theenumi}},leftmargin=*]
        \item Geschäftsordnung (60 Min.)
            \begin{enumerate}
                \item Fortgeltung der alten Geschäftsordnung (V1718-002)
                \item Änderungen an der Geschäftsordnung
                \item Diskussion über Notwendigkeit von GO-Reform und
                Satzungsänderungen (V1718-001, V1718-003)
            \end{enumerate}
        \item Wahl des StuPa-Präsidiums (60 Min.)
        \item Wahl des Satzungs- Wahlordnungs- und  Geschäftsordnungsausschusses (15 Min.)
        \item Wahl des Ausschusses gegen Rechts (15 Min.)
        \item Wahl des Haushaltsausschusses (15 Min.)
        \item Wahl des Wirtschaftsrats (15 Min.)
        \item Wahl des Ältestenrats (45 Min.)
            \begin{enumerate}
                \item Verfahren zur Wahl zum Ältestenrat (15 Min.)
                \item Wahl des Ältestenrats (30 Min.)
            \end{enumerate}
        \item RIS-Wahl (15 Min.)
            \begin{enumerate}
                \item Bestätigung der Wahlniederschrift
                \item Bestätigung der Referentinnen
            \end{enumerate}
        \item Queer-Wahl (15 Min.)
            \begin{enumerate}
                \item Bestätigung der Wahlniederschrift
                \item Bestätigung der Referentinnen
            \end{enumerate}
        \item RBCS-Wahl (15 Min.)
            \begin{enumerate}
                \item Bestätigung der Wahlniederschrift
                \item Bestätigung der Referentinnen
            \end{enumerate}
        \item Alle Frauen*-Wahl (15 Min.)
            \begin{enumerate}
                \item Bestätigung der Wahlniederschrift
                \item Bestätigung der Referentinnen
            \end{enumerate}
        \item Rechenschaftsbericht und Entlastung des amtierenden AStA (60 Min.)
            \begin{enumerate}
                \item Rechenschaftsbericht des amtierenden AStA
                \item Fragen und Diskussion
                \item Entlastung des AStA
              \end{enumerate}
        \item Wahl des neuen AStA-Vorstandes (60 Min.)
            \begin{enumerate}
                \item Diskussion VS-Thesen
                \item Wahl des AStA-Vorstandes
            \end{enumerate}
        \item Bestätigung der AStA-Referent*innen (90 Min.)
        \item Verschiedenes (10 Min.)
    \end{enumerate}

    \newpage

    % Hier beginnt Teil A. Bei Fortsetzungen kommt hier nur der Teil über
    % die Feststellung der endgültigen Fassung der Tagesordnung.
    % Für sonstige Sitzungen ist dies bereits das Gerüst für die Formalia.
    % Man muss nur noch die protokollierten Inhalte eintragen.

    {\Large \underline{Teil A}}

    \subsection{Formalia}

    \subsubsection{Geschäftsbericht Präsidium}

    \subsubsection{Anfragen an das Präsidium}

    Keine.

    \subsubsection{Geschäftsbericht AStA}


    \subsubsection{Anfragen an den AStA}


    \subsubsection{Dringlichkeitsanträge des AStA}

    Keine.

    \subsubsection{Aktuelle Stunde}

    \subsubsection{Feststellung der endgültigen Fassung des Teils B der Tagesordnung}
    \newpage
    %%% BESCHLOSSENE TAGESORDNUNG %%%
    \underline{Endgültige Fassung der Tagesordnung}

    \begin{enumerate}[label={\textbf{Top \theenumi}},leftmargin=*]
        \item Geschäftsordnung (60 Min.)
            \begin{enumerate}
                \item Fortgeltung der alten Geschäftsordnung (V1718-002)
                \item Änderungen an der Geschäftsordnung
                \item Diskussion über Notwendigkeit von GO-Reform und
                Satzungsänderungen (V1718-001, V1718-003)
            \end{enumerate}
        \item Wahl des StuPa-Präsidiums (60 Min.)
        \item Wahl des Satzungs- Wahlordnungs- und  Geschäftsordnungsausschusses (15 Min.)
        \item Wahl des Ausschusses gegen Rechts (15 Min.)
        \item Wahl des Haushaltsausschusses (15 Min.)
        \item Wahl des Wirtschaftsrats (15 Min.)
        \item Wahl des Ältestenrats (45 Min.)
            \begin{enumerate}
                \item Verfahren zur Wahl zum Ältestenrat (15 Min.)
                \item Wahl des Ältestenrats (30 Min.)
            \end{enumerate}
        \item RIS-Wahl (15 Min.)
            \begin{enumerate}
                \item Bestätigung der Wahlniederschrift
                \item Bestätigung der Referentinnen
            \end{enumerate}
        \item Queer-Wahl (15 Min.)
            \begin{enumerate}
                \item Bestätigung der Wahlniederschrift
                \item Bestätigung der Referentinnen
            \end{enumerate}
        \item RBCS-Wahl (15 Min.)
            \begin{enumerate}
                \item Bestätigung der Wahlniederschrift
                \item Bestätigung der Referentinnen
            \end{enumerate}
        \item Alle Frauen*-Wahl (15 Min.)
            \begin{enumerate}
                \item Bestätigung der Wahlniederschrift
                \item Bestätigung der Referentinnen
            \end{enumerate}
        \item Rechenschaftsbericht und Entlastung des amtierenden AStA (60 Min.)
            \begin{enumerate}
                \item Rechenschaftsbericht des amtierenden AStA
                \item Fragen und Diskussion
                \item Entlastung des AStA
              \end{enumerate}
        \item Wahl des neuen AStA-Vorstandes (60 Min.)
            \begin{enumerate}
                \item Diskussion VS-Thesen
                \item Wahl des AStA-Vorstandes
            \end{enumerate}
        \item Bestätigung der AStA-Referent*innen (90 Min.)
        \item Verschiedenes (10 Min.)
    \end{enumerate}

    \subsubsection{Feststellung der Beschlussfähigkeit}

    \subsubsection{Genehmigung der Protokolle der vorangegangenen Sitzungen}


    \vspace{0.5cm}
    {\Large \underline{Teil B}}

    \subsection{Geschäftsordnung}

        \textbf{a) Fortgeltung der alten Geschäftsordnung (V1718-002)}

        \textbf{b) Änderungen an der Geschäftsordnung}

        \textbf{c) Diskussion über Notwendigkeit von GO-Reform und
        Satzungsänderungen (V1718-001, V1718-003)}

    \subsection{Wahl des StuPa-Präsidiums}

    \subsection{Wahl des Satzungs- Wahlordnungs- und  Geschäftsordnungsausschusses}

    \subsection{Wahl des Ausschusses gegen Rechts}

    \subsection{Wahl des Haushaltsausschusses}

    \subsection{Wahl des Wirtschaftsrats}

    \subsection{Wahl des Ältestenrats}

        \textbf{a) Verfahren zur Wahl zum Ältestenrat}

        \textbf{b) Wahl des Ältestenrats}

    \subsection{RIS-Wahl}

        \textbf{a) Bestätigung der Wahlniederschrift}

        \textbf{b) Bestätigung der Referentinnen}

    \subsection{Queer-Wahl}

        \textbf{a) Bestätigung der Wahlniederschrift}

        \textbf{b) Bestätigung der Referentinnen}

    \subsection{RBCS-Wahl}

        \textbf{a) Bestätigung der Wahlniederschrift}

        \textbf{b) Bestätigung der Referentinnen}

    \subsection{Alle Frauen*-Wahl}

        \textbf{a) Bestätigung der Wahlniederschrift}

        \textbf{b) Bestätigung der Referentinnen}

    \newpage

    \subsection{Rechenschaftsbericht und Entlastung des amtierenden AStA}

        \textbf{a) Rechenschaftsbericht des amtierenden AStA}

        \textbf{b) Fragen und Diskussion}

        \textbf{c) Entlastung des AStA}

    \subsection{Wahl des neuen AStA-Vorstandes}

        \textbf{a) Diskussion VS-Thesen}

        \textbf{b) Wahl des AStA-Vorstandes}

    \subsection{Bestätigung der AStA-Referent*innen}

    \subsection{Verschiedenes}

\end{document}
